\documentclass[11pt]{article}
\usepackage[a4paper,margin=1in]{geometry}
\usepackage[T1]{fontenc}
\usepackage{lmodern}
\usepackage{microtype}
\usepackage{amsmath,amssymb,amsthm,mathtools}
\usepackage{booktabs}
\usepackage{xcolor}
\usepackage[numbers]{natbib}
\usepackage[colorlinks=true,allcolors=blue]{hyperref}

\newtheorem{theorem}{Theorem}
\newtheorem{corollary}[theorem]{Corollary}
\newcommand{\Tr}{\operatorname{Tr}}

\title{Block-Power Quotient Envelope Criterion}
\author{Bin Wang}
\date{July 2026}

\begin{document}
\maketitle

\begin{abstract}
The orthogonal quotient of RH-245 has no one-step contraction, so a direct
geometric estimate is unavailable.  We prove the exact block-power substitute:
a uniformly bounded trace norm and a uniformly contractive operator norm for
one power $C^m$ imply a geometric envelope for all later quotient traces and
an explicit logarithmic tail bound.  A finite $m=12$ diagnostic on 17
tractable endpoints gives rate $0.393300$ and unit-disk tail bound
$1.80\times10^{-5}$.  The constants are not promoted to a noise-uniform
certificate.
\end{abstract}

\section{Block criterion}

Let $C_\sigma\in\mathcal S_2(H)$ be the orthogonal quotient from RH-245 and
$\tau_n(\sigma)=\Tr C_\sigma^n$ for $n\ge2$ \citep{WangRH245}.  Fix an integer
$m\ge2$.  Suppose
\[
 K_m=\sup_\sigma\|C_\sigma^m\|_1<\infty,
 \qquad
 \eta_m=\sup_\sigma\|C_\sigma^m\|<1,
\]
and set $L_r=\sup_\sigma\|C_\sigma^r\|<\infty$ for
$0\le r<m$ (with $L_0=1$).

\begin{theorem}[block-power envelope]\label{thm:block}
For $n=\ell m+r\ge m$, $\ell\ge1$, $0\le r<m$,
\begin{equation}\label{eq:block-bound}
 |\tau_n(\sigma)|
 \le K_m L_r\eta_m^{\ell-1}.
\end{equation}
Writing $q=\eta_m^{1/m}$, the right side is at most $M_mq^n$, where
\begin{equation}\label{eq:M}
 M_m=K_m\max_{0\le r<m}L_rq^{-(m+r)}.
\end{equation}
If the finite head $2\le n<m$ has a uniform bound compatible with the same
rate, then $|\tau_n(\sigma)|\le Mq^n$ for every $n\ge2$ and RH-240 applies.
\end{theorem}

\begin{proof}
Factor $C_\sigma^n=(C_\sigma^m)^\ell C_\sigma^r$ and use
$|\Tr T|\le\|T\|_1$, submultiplicativity, and the ideal inequality
$\|UV\|_1\le\|U\|_1\|V\|$.  This gives \eqref{eq:block-bound}.  Since
$\eta_m^{\ell-1}=q^{n-(m+r)}$, \eqref{eq:M} follows.  The finite head is
handled by enlarging $M$ finitely; RH-240 then supplies local normality and
zero-freeness of the relative determinant \citep{WangRH240}.
\end{proof}

The same factorization yields a useful tail estimate.  For $R\ge0$ with
$\eta_mR^m<1$,
\begin{equation}\label{eq:tail}
 \sum_{n\ge m}\frac{|\tau_n(\sigma)|R^n}{n}
 \le
 \frac{K_mR^m}{m(1-\eta_mR^m)}
 \sum_{r=0}^{m-1}L_rR^r.
\end{equation}
This is a certificate template, not a claim that its hypotheses hold
uniformly for the moving cloud.

\section{Finite 12-block diagnostic}

We use the 17 ordered-Schur quotient rows of RH-245 (dimensions at most 512)
and take maxima over this finite subbatch.  The first contractive depth is
3--7, with no one-step contraction.  For $m=12$,
\[
 \eta_{12}\le1.3698766308677744\times10^{-5},\qquad
 K_{12}\le1.5684781027206807\times10^{-5},
\]
so
\[
 q_{12}=0.3932995547481413,
 \qquad
 \sum_{n\ge12}\frac{|\tau_n|}{n}
 \le1.7991531976413385\times10^{-5}
\]
under the finite-sample constants and $R=1$.  The corresponding finite
remainder operator-norm maxima for orders $0$ through $11$ are retained in
the JSON audit, and give $M_{12}=697.6520352674352$ through
\eqref{eq:M}.  The large $M$ is a reminder that a fast asymptotic block rate
does not by itself make the first uncontrolled orders small.

\begin{center}
\small
\begin{tabular}{lr}
\toprule
quantity & value\\
\midrule
finite endpoints & 17\\
first contractive depth & 3--7\\
one-step contractive endpoints & 0/17\\
$q_{12}$ & 0.3932995547\\
$M_{12}$ & 697.6520353\\
unit-disk tail from order 12 & $1.7992\times10^{-5}$\\
\bottomrule
\end{tabular}
\end{center}

\section{Boundary}

The theorem is an exact conditional route to the all-order envelope.  The
finite diagnostic does not establish uniform $K_m$, $\eta_m$, or $L_r$ over
all 32 endpoints, interval perturbations, or a continuum of noise.  It also
does not identify a cloud whose residual has the RH-243 deterministic anchor.
The separate-absolute route is tested next; it is expected to fail because
it discards the cancellation that makes the quotient small.

Gate A remains open and Gates B--E are untouched.  No Hilbert--P\'olya
operator, zeta-divisor equality, Riemann-zero identification, or RH
implication is claimed.

\bibliographystyle{plainnat}
\bibliography{references}
\end{document}
